% Bewertungspaket zum Gesamttest «Textaufgaben modellieren» — Quelle des PDFs (python3 scripts/build-lp-pdf.py modellieren).
% Ganze Punkte je Teilschritt; Werte und Fehlerbeispiele mit python3 nachgerechnet (scripts/lp/modellieren/zahlen.py, 08.10.2026).
\documentclass[11pt]{article}
\usepackage{../lp-druck}
\renewcommand{\lpname}{Textaufgaben modellieren}
\renewcommand{\lpbereich}{Grundlagenfach}
\renewcommand{\lpteilgebiet}{2.1 · 2.3}
\renewcommand{\lpfuss}{Bewertungspaket}
\newcolumntype{P}{>{\centering\arraybackslash}p{10mm}}
\newenvironment{raster}{\par\smallskip\noindent\tabularx{\linewidth}{@{}>{\raggedright\arraybackslash}p{38mm}P X@{}}\toprule
  \small\textsc{Teilschritt} & \small\textsc{P} & \small\textsc{Musterlösung}\\\midrule}{\bottomrule\endtabularx\par}
\newcommand{\fehler}[1]{\par\smallskip{\small\textcolor{lprot}{\bfseries Typische Fehler:} #1}}
\newcommand{\E}{\,{\footnotesize\bfseries\color{lpgruen}(E)}}
\begin{document}
\lpkopf{Bewertungspaket zum Gesamttest}{}

\begin{lpachtung}{Erst lösen, dann öffnen}
Dieses Paket enthält die Musterlösung. Wer es vor dem Lösen liest, misst am Ende nur, wie gut das Abschreiben gelingt.
\end{lpachtung}

\section*{1 · So gehst du vor}
\begin{enumerate}[leftmargin=*,itemsep=2pt]
\item \textbf{Gesamttest lösen} — vollständig, mit Deklaration, Ansatz und Rechenweg, ohne dieses Paket. Teil A ohne Taschenrechner.
\item \textbf{Fotografieren oder scannen}, gut lesbar, eine Seite pro Bild. \textbf{Kein Name} auf den Bildern.
  Handyfotos tragen oft unsichtbar Standort, Zeit und Gerät mit (Metadaten). Lade darum lieber einen Scan oder ein
  Bildschirmfoto des Fotos hoch, oder entferne vorher die Standortangaben.
\item \textbf{Bewerten lassen.} Ob du dafür eine KI nutzt und welche, entscheidest du selbst und in eigener Verantwortung —
  je nachdem, was dir aufgrund deines Alters und deiner persönlichen Situation erlaubt ist (Altersgrenzen und
  Nutzungsbedingungen des Anbieters, allenfalls Einverständnis deiner Eltern). Lade nur deine Lösung und dieses PDF hoch,
  keine persönlichen Daten, und füge den Auftrag aus Abschnitt~2 ein. Ohne KI kannst du dich mit Abschnitt~3 selbst bewerten.
\item \textbf{Rückmeldung prüfen.} Stimmt, was die KI aus deiner Handschrift gelesen hat? Eine KI kann sich irren —
  im Zweifel gilt das Raster in Abschnitt~3.
\item \textbf{Gezielt wiederholen} nach Abschnitt~4.
\end{enumerate}

\section*{2 · Auftrag an die KI}
\begin{tcolorbox}[colback=lplinie!25,colframe=lplinie,boxrule=0.5pt,arc=1mm,fontupper=\small]
Du bist Korrektorin oder Korrektor für Mathematik an einer Schweizer Berufsmaturitätsschule. Im Anhang findest du (1)~das Bewertungspaket zum Gesamttest «Textaufgaben modellieren» mit Musterlösung und Punkteraster und (2)~Fotos meiner handschriftlichen Lösung.\par\medskip
Geh so vor:\par
1. Schreib zu jeder Aufgabe G1 bis G7 zuerst ab, was du in meiner Lösung liest — Deklaration, Ansatz, Rechenweg und Ergebnis. Was du nicht sicher lesen kannst, markierst du mit [unsicher]. Nicht raten.\par
2. Vergib die Punkte genau nach dem Raster im Bewertungspaket (Abschnitt 3), ganze Punkte je Zeile. Ziehe einen Fehler nur einmal ab: Rechne ich mit einem falschen Zwischenergebnis richtig weiter, gibt es die Folgepunkte. Ausnahme: Zeilen mit \textbf{(E)} gibt es nur für das richtige Ergebnis, nie als Folgepunkt. Die Zeilen «Typische Fehler» sagen, welche Rasterzeile bei diesem Fehler 0 Punkte gibt — sie sind kein zusätzlicher Abzug. Führen zwei verschiedene Fehler zu denselben Zahlen, bewertest du sie gleich.\par
3. Andere richtige Lösungswege zählen voll: andere Namen für die Unbekannten, eine statt zwei Unbekannte (oder umgekehrt), Einsetzen, Gleichsetzen oder Addition, mit 100 oder 1000 multiplizierte Gleichungen, Dezimalpunkt oder Dezimalkomma, Bruch oder Dezimalzahl, Prozent oder Dezimalzahl in der Antwort. Zeilen mit \textbf{(E)} sind Ergebnispunkte: Das Ergebnis muss stimmen \emph{und} aus einem erkennbaren Weg hervorgehen. Alle übrigen Zeilen bewerten den Weg selbst.\par
4. Begründe jeden Abzug in einem Satz und nenne die Fehlerart (zum Beispiel «Prozent statt Stoffmenge» oder «Zeitanteil vergessen»). Schreib die Musterlösung nicht ab — zeig mir, wo mein Fehler liegt.\par
5. Gib zum Schluss eine Tabelle «Aufgabe | Punkte | Begründung», die Gesamtpunktzahl (von 25) und — nach der Selbsteinschätzung in Abschnitt 4 — welches Kapitel des Leitprogramms ich wiederholen soll. Keine Note.\par
6. Fehlt ein Foto oder ist eine Aufgabe unlesbar, sag es, statt sie zu bewerten.
\end{tcolorbox}

\needspace{16\baselineskip}
\section*{3 · Musterlösung und Punkteraster}
\textbf{Allgemein:} Ganze Punkte je Zeile. Ein Fehler wird nur einmal abgezogen (Folgefehler). Gleichwertige Wege und
Schreibweisen zählen voll. In der Spalte P: \textbf{(E)}~= Ergebnispunkt, Ergebnis richtig \emph{und} Weg erkennbar;
ohne Zeichen = die Zeile bewertet den Weg oder die Begründung. \textbf{Teil A} ist ohne Taschenrechner gestellt: Ein dort nur
mit dem Rechner gelöstes System (ohne Rechenweg von Hand) gibt die Lösungszeile nicht.
\textbf{Deklaration:} Sie zählt nur mit Bedeutung \emph{und} Einheit bzw. Bezug («\(x\): Masse der 10-\%-Lösung in kg»); «\(x\) = Lösung» gibt 0.
\textbf{Teilrichtige Zeilen:} Verlangt eine Zeile mehrere Angaben und stimmt nur ein Teil, gibt sie 0 Punkte.
\textbf{Fehlt eine Lösung} oder ist eine unbrauchbare nicht verworfen, gibt die Antwortzeile 0 — auch als Folgefehler.
\textbf{Gleiche Zahlen, gleiche Punkte:} Führen zwei verschiedene Fehler zu denselben Zahlen, nennt das Raster beide und bewertet sie gleich.
\textbf{Folgepunkte für Grundform und Eingabe} gibt es, wenn sie zum eigenen (falschen) Ansatz passen. Eine geordnete Form
\(a \cdot x + b \cdot y = c\) bzw. \(a \cdot x^2 + b \cdot x + c = 0\) gilt als notierte Eingabe, auch ohne «eingetippt: …».

\begin{lpaufgabe}{G1}{Ziffernrätsel}{4 P}
\begin{raster}
Deklaration & 1 & \(z\): Zehnerziffer, \(e\): Einerziffer (Ziffern, \(z \neq 0\)).\\
System & 1 & \(10 \cdot z + e = 4 \cdot (z + e) + 9\) und \(10 \cdot e + z = 10 \cdot z + e + 27\). Beide nötig.\\
Lösung von Hand & 1\E & Geordnet: \(6 \cdot z - 3 \cdot e = 9\), also \(2 \cdot z - e = 3\); und \(e - z = 3\). Addiert: \(z = 6\), \(e = 9\). Die Zahl heisst \textbf{69}.\\
Probe am Text & 1 & \(4 \cdot 15 + 9 = 69\) und \(96 - 69 = 27\). Gilt als Folgefehler, wenn die Probe zur eigenen Lösung richtig gerechnet ist und einen Fehler aufdeckt (so benannt).\\
\end{raster}
\fehler{«um 9 grösser» auf der falschen Seite (\(10 \cdot z + e + 9 = 4 \cdot (z + e)\)): Systemzeile 0; mit \(e - z = 3\) folgt \(z = 0\) — keine zweistellige Zahl; wer das bemerkt und so sagt, erhält die Probezeile.
Stellenwert vergessen (\(z + e\) statt \(10 \cdot z + e\)): Systemzeile 0, Lösungszeile 0 (E).
Zahl nur durch Probieren gefunden, ohne das System zu lösen: Lösungszeile 0 (verlangt ist «von Hand lösen»), die übrigen Zeilen nach dem Aufgeschriebenen.
Mit dem Taschenrechner gelöst ohne Rechenweg: Lösungszeile 0.}
\end{lpaufgabe}

\begin{lpaufgabe}{G2}{Ansätze vergleichen}{4 P}
\begin{raster}
(a) richtige Ansätze & 1 & B und C sind richtig (C mit einer Unbekannten, der Rest \(300 - x\) als Term). Beide nötig.\\
(a) Fehler in A & 1 & In A sind die rechten Seiten vertauscht: Die Stückbilanz zählt Plätze (300), die Wertbilanz Franken (7400).\\
(b) Lösung von Hand & 1\E & C: \(22 \cdot x + 9000 - 30 \cdot x = 7400\), \(-8 \cdot x = -1600\), \(x = 200\). Ebenso mit B (Einsetzen oder Addition).\\
(b) Antwort & 1 & \textbf{200 Parkett- und 100 Balkonplätze}. Folgepunkt, wenn zur eigenen Lösung beide Anzahlen genannt sind und beide ganz und nicht negativ sind.\\
\end{raster}
\fehler{Nur B als richtig erkannt (C verworfen, «es braucht zwei Gleichungen»): erste Zeile 0.
Preise vertauscht (\(30 \cdot x + 22 \cdot (300 - x) = 7400\)), gibt \(x = 100\): Lösungszeile 0; die Antwort «100 Parkett- und 200 Balkonplätze» erhält den Folgepunkt.
Minus vor der Klammer vergessen (\(22 \cdot x + 9000 + 30 \cdot x = 7400\)): Lösungszeile 0; \(x = -\tfrac{1600}{52}\) ist keine Anzahl — Antwortzeile 0, auch wenn das so gemeldet wird.}
\end{lpaufgabe}

\begin{lpaufgabe}{G3}{Menü und Einzelverkauf}{3 P}
\begin{raster}
Deklaration und Stückbilanzen & 1 & \(x\): einzeln verkaufte Sandwiches, \(y\): einzeln verkaufte Getränke, \(z\): Anzahl Menüs. \(x + z = 50\) (Sandwiches) und \(y + z = 70\) (Getränke): Das Menü zählt in beiden. Alles nötig.\\
Wertbilanz & 1 & \(6 \cdot x + 3 \cdot y + 8 \cdot z = 476\); gleichwertig mit \(x = 50 - z\), \(y = 70 - z\) eingesetzt: \(6 \cdot (50 - z) + 3 \cdot (70 - z) + 8 \cdot z = 476\).\\
Lösung von Hand & 1\E & \(300 + 210 - z = 476\), \(z = 34\): \textbf{34 Menüs} (dazu 16 Sandwiches und 36 Getränke einzeln). Probe: \(96 + 108 + 272 = 476\).\\
\end{raster}
\fehler{Menü nur in einer Stückbilanz gezählt (etwa \(y = 70\)): Zeile «Stückbilanzen» 0; es folgt \(z = -17\), keine Anzahl — Lösungszeile 0.
Menüpreis als Summe der Einzelpreise (9 CHF statt 8 CHF): Wertbilanzzeile 0, Lösungszeile 0.
Ein anderer richtiger Weg (Addition, Gleichungen subtrahiert, \(x\) und \(y\) zuerst bestimmt) zählt voll.}
\end{lpaufgabe}

\begin{lpaufgabe}{G4}{Zwei Lösungen und Wasser}{3 P}
\begin{raster}
Deklaration und System & 1 & \(x\): Masse der 10-\%-Lösung in kg, \(y\): Masse der 30-\%-Lösung in kg. \(x + y + 10 = 50\) und \(0.1 \cdot x + 0.3 \cdot y + 0 \cdot 10 = 0.14 \cdot 50 = 7\). Beide nötig.\\
Eingabe sys-solv & 1 & Zeile 1: 1, 1, 40; Zeile 2: 0.1, 0.3, 7 (oder mal 10: 1, 3, 70; mal 100: 10, 30, 700). Folgepunkt, wenn die Eingabe zum eigenen System passt.\\
Lösung & 1\E & \(x = 25\), \(y = 15\): \textbf{25 kg der 10-\%-Lösung und 15 kg der 30-\%-Lösung.} Probe: \(2.5 + 4.5 = 7\) kg Wirkstoff in 50 kg.\\
\end{raster}
\fehler{Wasser in der Mengenbilanz vergessen (\(x + y = 50\)): Systemzeile 0; gibt \(x = 40\), \(y = 10\) — Eingabe als Folgefehler 1, Lösungszeile 0 (E): 1 von 3 P.
Stoff aus 40 kg statt aus 50 kg (\(0.14 \cdot 40 = 5.6\)): Systemzeile 0; gibt \(x = 32\), \(y = 8\) — Eingabe 1, Lösungszeile 0: 1 von 3 P.
Rechts der Anteil 0.14 statt der Stoffmenge: Systemzeile 0, Lösungszeile 0 (negative Masse).}
\end{lpaufgabe}

\begin{lpaufgabe}{G5}{Eindampfen}{4 P}
\begin{raster}
Deklaration und System & 1 & \(m\): Masse der Lösung vor dem Verdunsten in kg, \(p\): Zuckeranteil vorher (Dezimalzahl). \(m \cdot p = 6\) und \((m - 10) \cdot (p + 0.05) = 6\). Alles nötig.\\
Grundform & 1 & Ausmultipliziert mit \(m \cdot p = 6\): \(0.05 \cdot m - 10 \cdot p = 0.5\), \(p = 0.005 \cdot m - 0.05\); eingesetzt und mal 200: \(m^2 - 10 \cdot m - 1200 = 0\) (oder ein Vielfaches). Gleichwertig über \(m = 200 \cdot p + 10\): \(200 \cdot p^2 + 10 \cdot p - 6 = 0\).\\
Lösung & 1\E & poly-solv: \(m = 40\) (bzw. \(p = 0.15\)); \(p = 6 : 40 = 0.15\): \textbf{40 kg Lösung mit 15\,\% Zucker}. Beide Angaben nötig. Probe: 30 kg mit 6 kg Zucker sind 20\,\%.\\
Verwerfen & 1 & \(m = -30\) (bzw. \(p = -0.2\)) fällt weg: Eine Masse (ein Anteil) ist nie negativ. Folgepunkt, wenn zur eigenen Grundform die negative Lösung begründet verworfen wird.\\
\end{raster}
\fehler{«5 Prozentpunkte» als \(p + 5\): Systemzeile 0; es folgt \(m^2 - 10 \cdot m - 12 = 0\), \(m \approx 11.08\) — Grundform als Folgefehler 1, Lösungszeile 0 (E), Verwerfen als Folgefehler 1.
\(m + 10\) statt \(m - 10\) (Wasser dazugegeben statt verdunstet): Systemzeile 0; die Grundform \(m^2 + 10 \cdot m + 1200 = 0\) hat keine Lösung — Lösungs- und Verwerfenzeile 0.
Vorzeichenfehler beim Einsetzen (\(m^2 + 10 \cdot m - 1200 = 0\), gibt \(m = 30\), \(p = 0.2\)): Grundformzeile 0, Lösungszeile 0 (E), Verwerfen als Folgefehler 1; die Probe (20 kg mit 25\,\% sind 5 kg Zucker, nicht 6) deckt den Fehler auf.
Nur \(m = 40\) ohne den Zuckeranteil: Lösungszeile 0.}
\end{lpaufgabe}

\begin{lpaufgabe}{G6}{Zinseszins mit Abhebung}{4 P}
\begin{raster}
Ansatz & 1 & \(p\): Zinssatz als Dezimalzahl. \((6000 \cdot (1 + p) - 1000) \cdot (1 + p) = 5110.60\), gleichwertig \(6000 \cdot (1 + p)^2 - 1000 \cdot (1 + p) = 5110.60\).\\
Grundform & 1 & \(6000 \cdot p^2 + 11\,000 \cdot p - 110.6 = 0\) oder ein Vielfaches (mal 5: \(30\,000 \cdot p^2 + 55\,000 \cdot p - 553 = 0\)). Gleichwertig mit \(q = 1 + p\): \(6000 \cdot q^2 - 1000 \cdot q - 5110.6 = 0\); mit \(p\) in Prozent: \(0.6 \cdot p^2 + 110 \cdot p - 110.6 = 0\).\\
Lösung & 1\E & \(p = 0.01\): \textbf{Zinssatz 1\,\%} (über \(q\): \(q = 1.01\); in Prozent: \(p = 1\)). Ob der Rechner \(\tfrac{1}{100}\) oder \(0.01\) zeigt, ist gleichgültig. Probe: \(6060 - 1000 = 5060\), \(5060 \cdot 1.01 = 5110.60\).\\
Verwerfen & 1 & Die zweite Lösung \(p = -\tfrac{553}{300} \approx -1.843\) (bzw. \(q \approx -0.843\), \(p \approx -184.3\) in Prozent) fällt weg: ein Zinssatz unter \(-100\,\%\) ist unmöglich. Gilt auch, wenn nur begründet wird, dass ein negativer Zinssatz hier nicht passt.\\
\end{raster}
\fehler{\(b = 12\,000\) statt \(11\,000\) — entsteht auf zwei Wegen, beide gleich bewertet: (1) die 1000 CHF nicht mitverzinst (\(6000 \cdot (1 + p)^2 - 1000 = 5110.60\)): Ansatzzeile 0, Grundform als Folgefehler 1; (2) richtiger Ansatz, aber \(-1000 \cdot p\) beim Ausmultiplizieren verloren: Ansatzzeile 1, Grundformzeile 0. In beiden Fällen \(p \approx 0.0092\): Lösungszeile 0 (E), Verwerfen als Folgefehler 1 \(\to\) 2 von 4\,P.
Mittelglied \(2 \cdot 6000 \cdot p\) von \((1 + p)^2\) vergessen (\(6000 \cdot p^2 - 1000 \cdot p - 110.6 = 0\), gibt \(p \approx 0.243\)): Grundformzeile 0, Lösungszeile 0 (E), Verwerfen als Folgefehler 1.
Beide Lösungen als Antwort: Verwerfen-Zeile 0.}
\end{lpaufgabe}

\begin{lpaufgabe}{G7}{Gleich viel Zins}{3 P}
\begin{raster}
Deklaration und System & 1 & \(x\): Kapital auf Anlage A in CHF, \(y\): Kapital auf Anlage B in CHF. \(x + y = 18\,000\) und \(0.02 \cdot \tfrac{9}{12} \cdot x = 0.012 \cdot y\), also \(0.015 \cdot x - 0.012 \cdot y = 0\). Zeitanteil nötig.\\
Eingabe sys-solv & 1 & Zeile 1: 1, 1, 18\,000; Zeile 2: 0.015, \(-0.012\), 0 (oder 15, \(-12\), 0; 5, \(-4\), 0). Folgepunkt, wenn die Eingabe zum eigenen System passt. Mit einer Unbekannten von Hand gelöst (\(0.027 \cdot x = 216\)) zählt die geordnete Gleichung.\\
Lösung & 1\E & \(x = 8000\), \(y = 10\,000\): \textbf{A 8000 CHF, B 10\,000 CHF, jede bringt 120 CHF Zins.} Alle drei Angaben nötig.\\
\end{raster}
\fehler{Zeitanteil vergessen (\(0.02 \cdot x = 0.012 \cdot y\)): Systemzeile 0; gibt \(x = 6750\), \(y = 11\,250\), je 135 CHF — Eingabe als Folgefehler 1, Lösungszeile 0 (E): 1 von 3 P.
9 statt \(\tfrac{9}{12}\) (\(0.18 \cdot x = 0.012 \cdot y\), gibt \(x = 1125\)): ebenso 1 von 3 P.
Zinssätze vertauscht zugeordnet (\(0.012 \cdot x = 0.015 \cdot y\), gibt \(x = 10\,000\), \(y = 8000\)): ebenso 1 von 3 P.
Zinssätze in Prozent (\(2 \cdot \tfrac{3}{4} \cdot x = 1.2 \cdot y\)): Das ist die Zinsgleichung mal 100 — gleichwertig, zählt voll (anders als bei einer Gleichung mit einem Zinsbetrag rechts).}
\end{lpaufgabe}

\section*{4 · Selbsteinschätzung}
\begin{tabularx}{\linewidth}{@{}l X@{}}\toprule
22 – 25 P & Die geprüften Aufgabenarten sitzen. Wo du Punkte verloren hast: das Kapitel dieser Aufgabe nochmals (Zuordnung unten).\\
17 – 21 P & Den schwächsten Teil nochmals: Tüfteln, Kontrollclips und Übungen des Kapitels, in dem du die meisten Punkte verloren hast.\\
11 – 16 P & Zurück zu den Kapiteln aller Aufgaben, in denen du Punkte verloren hast — zuerst die Kontrollclips.\\
0 – 10 P & Zurück zu Kapitel 0 und von dort der Reihe nach weiter.\\\midrule
\multicolumn{2}{@{}p{\linewidth}@{}}{\small\textbf{Aufgabe $\to$ Kapitel:} G1 $\to$ Kapitel 1 (Zahlen- und Ziffernrätsel); G2, G3 $\to$ Kapitel 3 (Verteilen); G4, G5 $\to$ Kapitel 2 (Mischen); G6, G7 $\to$ Kapitel 4 (Zins)}\\\bottomrule
\end{tabularx}
\end{document}
